\section{Análise e Redesenho do Processo}

Esta seção materializa a etapa de \textit{Redesenho} do ciclo
de BPM apresentado na Base Teórica, articulando o diagnóstico
do processo atual (AS-IS), descrito na seção de Levantamento
de Requisitos, com a especificação do processo futuro
(TO-BE), que orienta o desenvolvimento da solução apresentado
na seção seguinte. O objetivo não é apenas contrastar duas
representações do fluxo, mas evidenciar, campo a campo,
decisão a decisão, a origem de cada mudança nos requisitos
levantados junto aos três perfis consultados.

\subsection{Processo Atual (AS-IS)}

O diagnóstico consolidado a partir das respostas da Analista,
da Supervisora e do Gerente caracteriza um processo
descentralizado e dependente de controles manuais. Os
chamados internos eram registrados no SAP Service --- mesma
ferramenta utilizada como CRM para atendimentos e reclamações
de clientes externos ---, misturando ocorrências internas e
manifestações externas em um único ambiente. A chegada do
chamado era comunicada paralelamente por um grupo de
WhatsApp, sem protocolo automático, sem controle de prazo e
sem visibilidade gerencial sobre volume, prioridade ou tempo
de resolução. O Quadro~\ref{quadro:as-is} sintetiza as
características centrais do processo atual.

\begin{table}[H]
\caption{Características do processo atual (AS-IS)}
\label{quadro:as-is}
\centering
\begin{tabular}{p{3.6cm} p{8.2cm}}
\hline
\textbf{Dimensão} & \textbf{Situação observada} \\
\hline
Canal de abertura & Cadastro ``Experiência do Cliente'' no
SAP Service, compartilhado com manifestações externas de
clientes. \\
\hline
Comunicação & Complementar e manual, por grupo de WhatsApp,
sem registro formal do repasse. \\
\hline
Controle de prazo & Inexistente; sem SLA definido ou
monitorado. \\
\hline
Rastreabilidade & Baixa; mistura entre chamados internos e
externos dificulta a leitura de volumes. \\
\hline
Responsabilidade & Informal; dependência de retorno manual
das áreas responsáveis. \\
\hline
Indicadores & Inexistentes; ausência de base para gestão
baseada em dados. \\
\hline
\end{tabular}
\raggedright
\par
{\fontsize{10}{12}\selectfont Fonte: Elaborado pelo autor com
base na pesquisa de campo}
\end{table}

\subsection{Processo Futuro (TO-BE)}

O processo redesenhado responde a cada uma das fragilidades
identificadas no AS-IS, substituindo controles manuais por
regras automatizadas e centralizando o registro em uma única
fonte de dados. O Quadro~\ref{quadro:to-be} apresenta a
transformação proposta, dimensão a dimensão, em correspondência
direta com o Quadro~\ref{quadro:as-is}.

\begin{table}[H]
\caption{Transformação do processo: AS-IS para TO-BE}
\label{quadro:to-be}
\centering
\begin{tabular}{p{3.0cm} p{4.9cm} p{4.9cm}}
\hline
\textbf{Dimensão} & \textbf{AS-IS} & \textbf{TO-BE} \\
\hline
Canal de abertura & Cadastro compartilhado no SAP Service &
Site no SharePoint com aplicativo Power Apps próprio para
chamados internos \\
\hline
Comunicação & Grupo de WhatsApp, sem registro & Notificações
automáticas por e-mail, vinculadas a eventos da Microsoft
List \\
\hline
Controle de prazo & Inexistente & SLA por tipo e prioridade,
com contagem automatizada \\
\hline
Rastreabilidade & Baixa, dados misturados & Protocolo único,
histórico centralizado na Microsoft List \\
\hline
Responsabilidade & Informal & Classificação e atribuição
explícita pela Qualidade/SAC, com transferência rastreável \\
\hline
Indicadores & Inexistentes & Painel gerencial no Power BI
(volume, tempo médio, cumprimento de SLA) \\
\hline
\end{tabular}
\raggedright
\par
{\fontsize{10}{12}\selectfont Fonte: Elaborado pelo autor}
\end{table}

\subsection{Decisões de Redesenho}

A transição entre os dois estados do processo não decorreu da
transposição direta do fluxo manual para um ambiente digital,
mas de decisões de redesenho fundamentadas nos requisitos
levantados. Duas decisões merecem destaque.

A primeira separa o registro de chamados internos das
manifestações externas de clientes, respondendo diretamente à
fragilidade central identificada pela Analista no processo
AS-IS. Essa separação é também uma decisão de Engenharia de
Processos: ao isolar o fluxo interno em um canal próprio, a
organização ganha capacidade de leitura de volumes e de
governança que o modelo anterior não permitia.

A segunda desloca a definição de área responsável, prioridade,
tipologia e responsável pela tratativa do formulário inicial
de abertura para uma etapa intermediária de classificação,
sob responsabilidade da Qualidade/SAC. Essa decisão resulta da
convergência entre a necessidade de filtro interno apontada
pela Analista e a necessidade de transferência de responsável
apontada pela Supervisora, e ilustra como o redesenho do
processo (TO-BE) é uma consequência direta do levantamento de
requisitos, e não uma escolha arbitrária de arquitetura de
software.

Por fim, foi estabelecida uma etapa de retorno obrigatório à
Qualidade/SAC após o preenchimento da causa raiz e da ação
corretiva pela área responsável, antes do encerramento
definitivo do chamado. Essa regra opera como mecanismo de
validação previsto pelo ciclo de BPM, impedindo que a própria
área executora encerre um chamado sem revisão, e reforça a
rastreabilidade e a governança como objetivos centrais do
redesenho.

O processo futuro (TO-BE) especificado nesta seção é a base
direta para os componentes desenvolvidos e descritos na seção
de Desenvolvimento da Solução, na qual cada estado, campo e
regra de negócio aqui definidos é implementado no ecossistema
Microsoft 365.
